\scanpage{10}

1) $X$ espace topologique, $U,V$ ouverts, condts équivalentes
\[
\begin{aligned}
\text{a)}\quad & U=\complement\overline V,\quad V=\complement\overline U\\
\text{b)}\quad & U\cap V=\phi,\quad\complement(U\cup V)\subset\overline U\cap\overline V
 \qquad\text{(d'où ég.lté.)}\\
\text{c)}\quad & U\cap V=\phi,\quad X=\overline{U\cup V},\quad\dot U=\dot V\\
\text{d)}\quad & \overline U\cap\overline V=\dot U=\dot V.\ \inserted{\unsure{\Leftarrow|\leftarrow|?}}
\end{aligned}
\]

$U$ et $V$ se déterminent \unsure{mutuellem$^{\mathrm t}$|mutuellem}. Pour $U$ fixé, il existe $V$ ayant les propriétés précédentes ssi $V=\complement\overline U$ est dense dans $\complement U$.

\unsure{2|3}) $X$ espace topologique, $K\subset X$ fermé, $U_K^{\deleted{\unsure{\prime|2|3}}}$ et $\deleted{\unsure{V}}U'_K$ deux ouverts \unsure{rel$^{\mathrm{ts}}$|rel$^{\mathrm t}$.|relatifs} de $X$ tels que
\[
\resizebox{.99\linewidth}{!}{$\displaystyle
\overline U_K^{\deleted{\unsure{\prime|-|2}}}=K\smallsetminus\deleted{\unsure{V}}U'_K,\quad
\overline{\deleted{\unsure{V}}U'_K}=K\smallsetminus U_K,\quad
C=\overline U_K\cap\overline{U'_K}=
\underset{\deleted{\overline U_K\smallsetminus U_K\unsure{=}\unsure{\overline{V_K}|\overline{U'_K}}\smallsetminus\unsure{V_K|U'_K}}}
 {\overset{\inserted{\overline U_K\smallsetminus U_K=\overline{\deleted{\unsure{V|U}}U'_K}\smallsetminus\deleted{\unsure{V|U}}U'_K}}
 {\text{frontière commune}}}$}
\]
de $U_K$ et $U'_K$ \uline{dans $K$} \inserted{(fermé dans $X$)}. Soit $B=U_K\cap\dot K$, \deleted{donc} donc $B$ ouvert ds $\dot K$ et fermé dans $U_K$ donc \unsure{rel$^{\mathrm t}$.} fermé ds $X$; Alors on a, posant \inserted{si $B\cap C=\phi$.}
\[
\resizebox{\linewidth}{!}{$\displaystyle\left\{\begin{aligned}
&U=U_K\smallsetminus\deleted{\unsure{D|\overline U|B}}B_{\deleted{\unsure{K}}}=\operatorname{Int}(U_K)\\
&V=X\smallsetminus\overline U_K^{\deleted{\prime}}=
 (\underbrace{X\smallsetminus U_K}_{\inserted{X_1}})\smallsetminus C=\operatorname{Int}(X_1)\\
&U=\complement\overline V,\ V=\complement\overline U,\qquad
 \dot U=\dot V=\overline U\cap\overline V=\complement(U\cup V)=B\cup C\\
&\deleted{U\cap K=U\subset K,\ V\cap K=V_K}\quad
 \underbrace{V\subset X_1}_{\inserted{\text{diff. }C}}
 \underbrace{{}\subset\overline V}_{\inserted{\text{diff }B}}
\end{aligned}\right.$}
\]

\uline{Cor} \deleted{Si $B\cap C=\phi$ i.e. $U_K\smallsetminus U$\unsure{.|,} Supposons} \inserted{Pour que} \deleted{$B$} $X_1=\complement U_K$ soit loc fermé, il f. et s. que $B=U_K\cap\dot K$ soit \deleted{fermé dans $X$, \unsure{i.e.|alors}} \deleted{\inserted{dans $\unsure{B_0|B\cup}$}} \deleted{$B$ est ouvert et} fermé dans $X$, \deleted{\unsure{i|il}} ou encore que ce soit un sommand \unsure{topolog} \deleted{\unsure{ds|dans} $\dot K$, et \unsure{posant} $\unsure{\dot K|K}=B'\amalg B$} \deleted{\unsure{ds|dans} $\dot K$. Posant alors $\dot K=B'\amalg \unsure{B|B'}$} ou encore qu'il soit fermé ds $B\cup C$, ou encore que $\overline B\cap C=\phi$. [Il suffit pour ceci que $\dot K\cap C=\phi$ i.e. que $U_K\unsure{\cup|\amalg}V_K\supset K^{\bullet}$]. \deleted{i.e. $U_K$ et $V_K$ sont}
% archive sheet number: 9
